% common-scrbook/preamble/environments.tex
%
% Block environments for the scrbook target — workingnotes, callout, and
% pass-through stubs for the kaobook target's wrappers (segmentepigraph,
% segmentwidesection) so the kramdown converter is target-agnostic.

%----------------------------------------------------------------------------
% mdframed — used by workingnotes and callout. Loaded with the framemethod
% set to default (TikZ would give richer borders but isn't needed for the
% restrained register the scrbook target wants).
%----------------------------------------------------------------------------
\usepackage[framemethod=default]{mdframed}

%----------------------------------------------------------------------------
% Working Notes — review-mode-only environment for the trailing
% ## Working Notes section of each segment. Stage-1 ingest already strips
% the section in --public, so no LaTeX-side conditional is needed here.
%
% Visual register: thin rule above and below, subtle warm-gray foreground,
% \small body. Reads as set-aside workshop content — unmistakably distinct
% from finished argument, but not so visually heavy that it dominates the
% segment when it's present.
%----------------------------------------------------------------------------
\newmdenv[
  topline       = true,
  bottomline    = true,
  rightline     = false,
  leftline      = false,
  linewidth     = 0.4pt,
  linecolor     = wn-rule,
  innertopmargin    = 8pt,
  innerbottommargin = 8pt,
  innerleftmargin   = 0pt,
  innerrightmargin  = 0pt,
  skipabove     = \medskipamount,
  skipbelow     = \medskipamount,
  fontcolor     = wn-fg,
  % font option drives mdframed to apply the size/weight to the body via
  % its own group, which propagates across paragraph breaks and survives
  % the page-split that kills inline `\small` near the env's begin.
  font          = \small,
]{workingnotesframe}

\newenvironment{workingnotes}{%
  \begin{workingnotesframe}%
  \noindent\textit{\textbf{Working Notes.}}\par\smallskip%
}{%
  \end{workingnotesframe}%
}

%----------------------------------------------------------------------------
% Callouts — Obsidian-style `> [!type] title` blocks. Type vocabulary:
% warning / note / info / tip / danger. For the scrbook target we render
% all of them with the same restrained left-bar treatment, distinguished
% only by the type label in bold uppercase. Bright per-type tints would
% compete with the body register; the visual hierarchy stays quiet.
%----------------------------------------------------------------------------
\newmdenv[
  topline       = false,
  bottomline    = false,
  rightline     = false,
  leftline      = true,
  linewidth     = 2pt,
  linecolor     = callout-rule,
  innertopmargin    = 4pt,
  innerbottommargin = 4pt,
  innerleftmargin   = 10pt,
  innerrightmargin  = 0pt,
  skipabove     = \medskipamount,
  skipbelow     = \medskipamount,
]{calloutframe}

\NewDocumentEnvironment{callout}{m m}{%
  \begin{calloutframe}%
  \small%
  \noindent\textbf{\MakeUppercase{#1}}%
  \IfStrEq{#2}{}{}{\quad #2}%
  \par\smallskip%
}{%
  \end{calloutframe}%
}

%----------------------------------------------------------------------------
% Pass-through stubs for kaobook-target wrappers. The base AsfLatex
% converter (segment-mode) emits these; AsfScrbookLatex's volume-mode
% override doesn't, but a future segment-mode call (e.g., for a snippet
% rendered with render_fragment) might. Defined as identity so they don't
% break compilation if invoked.
%----------------------------------------------------------------------------
\NewDocumentEnvironment{segmentepigraph}{}{%
  \par\medskip%
  \begingroup%
  \small\itshape%
}{%
  \par%
  \endgroup%
  \par\medskip%
}

\NewDocumentEnvironment{segmentwidesection}{}{%
}{%
}
